\section{Related Work}
\label{sec:related_work}

\subsection{Multimodal Object Tracking}
Modern visual object tracking has advanced significantly by moving beyond single-modal (e.g., RGB) inputs to incorporate complementary sensors like Thermal, Depth, or Event cameras, mitigating challenges such as occlusion and low illumination~\cite{wang2023visevent}. Standard approaches rely on Siamese networks or Transformers~\cite{ostrack, seqtrack} to fuse multi-sensor streams. Recently, state-space models (SSMs) such as Mamba~\cite{zhu2024vision, liu2024vmamba} have also been introduced to tracking to model long-range temporal dependencies efficiently. While multimodal trackers~\cite{vipt, hou2024sdstrack, untrack} achieve high accuracy, they fundamentally assume that both modalities are fully and synchronously available at every frame. When applied in real-world scenarios where sensors frequently de-synchronize or drop out, their rigid fusion architectures fail, feeding noise into the tracking pipeline~\cite{ipt}. Our work breaks this assumption by designing a unified architecture, utilizing a Fast-iTPN encoder and a Mamba interaction neck, that maintains robust tracking performance even when modalities are intermittently missing.

\subsection{Multi-task Learning}

% general MTL
[To be modified to adapt to this paper (copied from my CVPR)]

Multi-task Learning (MTL)~\cite{caruana1997multitask} aims at learning a single network that jointly estimates accurate predictions for multiple desired tasks. Recent research on multi-task learning in computer vision can be broadly divided into two categories~\cite{vandenhende2021multi,ruder2017overview}. 
The first group aims at addressing the unbalanced optimization issues - each task's loss function often exhibits distinct scales and convergence behaviors, jointly minimizing them can lead to optimization conflicts and performance degradation. To address this issue, prior work proposes to estimate loss weights dynamically~\cite{kendall2018multi,liu2019end,guo2018dynamic,chen2018gradnorm,lin2019pareto,sener2018multi,liu2021towards}, mitigating conflicts between gradient conflicts~\cite{yu2020gradient,liu2021conflict,chen2020just,chennupati2019multinet++,suteu2019regularizing}, or aligning features with single-task models~\cite{li2020knowledge,li2022universal}. 

% ~\cite{liu2019end,vandenhende2020mti,ye2022inverted}

Our work is more related to the second one which aims at designing architectures~\cite{kokkinos2017ubernet,ruder2019latent,vandenhende2019branched,liang2018evolutionary,bragman2019stochastic,strezoski2019many,xu2018pad,zhang2019pattern,bruggemann2021exploring,bilen2016integrated,zhang2018joint} that better share information across tasks 
% . 
% Instead of sharing the entire encoder across all tasks, the architecture based methods focus on improving the network architecture 
using cross-task attention mechanisms~\cite{misra2016cross}, task-specific attention modules~\cite{liu2019end,bhattacharjee2023vision}, cross-task feature interaction~\cite{ye2022inverted,vandenhende2020mti}, gating strategies or {mixture of experts modules~\cite{bruggemann2020automated,guo2020learning,chen2023mod,fan2022m3vit,ye2023taskexpert}, visual prompting~\cite{ye2022taskprompter,liu2023hierarchical}} and distillation of multiple visual foundation models~\cite{ranzinger2024radio,heinrich2025radiov2,lu2024swiss}. However, these methods mainly capture cross-task relations within the 2D space, and two recent works~\cite {li2023multi,zheng2023multi} show that 3D-awareness is vital for learning more structured features that are shared and beneficial for all tasks by using NeRF as a decoder or a 3D-aware regularizer. However, MuvieNeRF~\cite{zheng2023multi} requires multiple views and ground-truth camera parameters, limiting its use for practical scenarios. While 3DMTL~\cite{li2023multi} does not require camera parameters during inference and can be applied to standard MTL settings with single-view input, it is shown that the method has limited capacity of leveraging multi-view data for enhancing the 3D-awareness. 











\subsection{Learning with Missing Modalities}
Handling missing or incomplete modalities is a long-standing challenge in robust multimodal learning~\cite{missing_survey, missing_1}. Prior strategies broadly fall into three categories. The first relies on generative imputation via GANs~\cite{gans} or Diffusion Models~\cite{croitoru2023diffusion} to synthesize missing data, which incurs prohibitive computational overhead for real-time tracking. The second category employs prompt-based learning or modality masking~\cite{missing_reconstraction} to teach models to ignore the missing streams. The third utilizes cross-modal distillation~\cite{missing_disstilation}, forcing a degraded student to mimic a complete-modality external teacher. While effective, these methods fundamentally treat missing modalities by \emph{passively coping} (e.g., masking them away) or rely on rigid external teacher models. In contrast, our approach shifts to \emph{active compensation}. Through the Bilateral Modality-specific feature Reconstruction and Hallucination (BMR) module, we use lightweight cross-modal projectors to actively hallucinate the absent features. Furthermore, our Curriculum Missing Augmentation (CMA) bypasses the need for an external teacher by employing online self-distillation, binding complete and degraded network passes dynamically as the missing rate is annealed over training.

\subsection{Mixture-of-Experts in Multimodal Fusion}
Mixture-of-Experts (MoE) models scale capacity efficiently by dynamically routing inputs to specialized expert networks. The success of large-scale MoE in language models~\cite{dai2024deepseekmoe, llama-moe} has inspired its application in multimodal vision tasks, where models like FuseMoE~\cite{han2025fusemoe} and Flex-MoE~\cite{yunflex-moe} route different sensory inputs to adapt to dynamic scenarios. However, existing multimodal MoE architectures typically enforce a uniform design for expert capacities and route primarily based on semantic content. In the context of sensor dropout, a uniform capacity wastes computation when a modality vanishes. Building upon our conference work~\cite{tan2025flextrack}, we address this by employing a Heterogeneous Mixture-of-Experts (HMoE) that scales expert capacities. In this paper, we tightly couple HMoE with the BMR module. This BMR-HMoE fusion not only adapts its computational footprint to the degree of modality availability but also intelligently allocates resources to fill in the missing feature gaps, an capability absent in standard MoE tracking designs.
